% Using InitechBase, Chapter 1: Overview (annex Sec 8.2 U1).  Sources:
% spec/samir/dbf_format.h and os/samir/fs/dbf.c (field types and widths,
% dbf_create), fs/dbt.c, fs/ndx.c, cmd/report.c (.FRM), samir_main.c (DO
% .PRG), the Limitations table generated from the engine
% (harness/pubs/samir_tables.c); the sample files pubs/examples/samir/disk.
\chapter{Overview}
\Gates{test-pubs-build, test-pubs-examples, test-dbf-fields, test-samir-repl}
This chapter describes the information InitechBase keeps, the form in which it
keeps it, and the files it reads and writes.

\section{Databases, Records and Fields}
A \term{database} is a collection of information about one kind of thing:
the employees of a company, its customers, or the invoices it has sent out.
InitechBase keeps each database in a \term{database file}, with the extension
\file{.DBF}.

A database is made up of \term[record]{records}, one for each employee, customer or
invoice.  Every record of a database has the same \term[field]{fields}: an employee
record of \file{STAFF.DBF} has a field for the employee number, one for the
last name, one for the first name, and so on.  The list of the fields, with
the name, type and width of each, is the \term{structure} of the database.

InitechBase numbers the records of a database 1, 2, 3 and so on, in the order
in which they were added.  The \term{record number} of a record does not
change unless \cmd{PACK} removes records before it.

The structure of the staff database of the Sample Programs and Utilities
diskette is:

\noindent\begin{tabular}{@{}llrrl@{}}
\hline
\textsf{\textbf{Field}} & \textsf{\textbf{Type}} & \textsf{\textbf{Width}} &
  \textsf{\textbf{Dec}} & \textsf{\textbf{Contents}}\\
\hline
\texttt{EMPNO} & Character & 4 & & Employee number\\
\texttt{LASTNAME} & Character & 14 & & Last name\\
\texttt{FIRSTNAME} & Character & 10 & & First name\\
\texttt{DEPT} & Character & 4 & & Department code\\
\texttt{EXT} & Character & 4 & & Telephone extension\\
\texttt{HIRED} & Date & 8 & & Date of hire\\
\texttt{SALARY} & Numeric & 9 & 2 & Annual salary\\
\texttt{EXEMPT} & Logical & 1 & & Exempt from overtime\\
\texttt{NOTES} & Memo & 10 & & Remarks\\
\hline
\end{tabular}

\section{Field Types}
Each field has one of five types.

\SideHead{Character}Letters, digits, blanks and any other characters, up to
254 of them.  A character field is always as wide as its structure says: a
shorter value is filled out with blanks on the right.

\SideHead{Numeric}A number, positive or negative, with a fixed number of
decimal places.  The width counts every digit, the decimal point and the
minus sign.  \texttt{SALARY}, of width 9 with 2 decimal places, holds numbers
up to 999999.99.  A number too large for its field is stored as asterisks.

\SideHead{Date}A date, stored as year, month and day.  A date field is 8
characters wide.  InitechBase displays dates in the form chosen by SET DATE,
with two digits for the year unless SET CENTURY is ON.  A date field may be
blank.  Dates can be compared and subtracted: the difference between two
dates is a number of days.

\SideHead{Logical}True or false, displayed \texttt{T} or \texttt{F} by
\cmd{LIST} and \texttt{.T.} or \texttt{.F.} by \cmd{?}.  A logical field is one
character wide.

\SideHead{Memo}Text of any length, such as notes or remarks.  The text is kept
in a separate \term{memo file} with the name of the database and the extension
\file{.DBT}; the record holds only the place of the text in the memo file.  A
memo field takes ten characters of the record.

\section{Files Used by InitechBase}
\noindent\hspace*{-\imHang}\begin{tabular}{@{}p{24mm}p{110mm}@{}}
\hline
\textsf{\textbf{Extension}} & \textsf{\textbf{File}}\\
\hline
\file{.DBF} & Database file: the structure and the records.  Opened with
  \cmd{USE}.\\
\file{.DBT} & Memo file: the text of the memo fields of the database of the same
  name.  Opened with the database.\\
\file{.NDX} & Index file: the records of a database in the order of a key, such as
  the last name.  Opened with \cmd{USE} \ldots{} \kw{INDEX}.\\
\file{.PRG} & Program file: InitechBase commands, run with \cmd{DO}.\\
\file{.FRM} & Report form file: the layout of a report, used by \cmd{REPORT
  FORM}.\\
\file{.TXT} & Text file: a report written by \cmd{REPORT FORM} \ldots{} \kw{TO
  FILE}.\\
\hline
\end{tabular}

Appendix~C describes the layout of the database, memo and index files.

\section{Specifications}
InitechBase works within the limits below.  They are given again at the
beginning of Chapter~5.

\noindent\hspace*{-\imHang}\begin{tabular}{@{}p{90mm}p{44mm}@{}}
\hline
\textsf{\textbf{Item}} & \textsf{\textbf{Limitation}}\\
\hline
\TableRows{samir_limits}
\hline
\end{tabular}

\section{The Dot Prompt}
When you start InitechBase, it displays a period and a blank, the \term{dot
prompt}, and waits for a command.  You type the command and press
\key{Enter}; InitechBase carries it out, displays its results and the dot
prompt again.  Chapter~2 describes how commands are typed.

\section{Commands by Task}
The commands of InitechBase fall into a few groups.  Chapter~5 describes each
command; the table names the commands of each task.

\noindent\hspace*{-\imHang}\begin{tabular}{@{}p{44mm}p{90mm}@{}}
\hline
\textsf{\textbf{Task}} & \textsf{\textbf{Commands}}\\
\hline
Opening and closing files & \cmd{USE}, \cmd{SELECT}, \cmd{CLOSE}, \cmd{QUIT}\\
Looking at records & \cmd{LIST}, \cmd{DISPLAY}, \cmd{?}, \cmd{??}\\
Moving the record pointer & \cmd{GO}, \cmd{SKIP}\\
Finding records & \cmd{LOCATE}, \cmd{CONTINUE}, \cmd{SEEK}, \cmd{FIND}\\
Adding and changing records & \cmd{APPEND BLANK}, \cmd{REPLACE}\\
Deleting records & \cmd{DELETE}, \cmd{RECALL}, \cmd{PACK}, \cmd{ZAP}\\
Reports & \cmd{REPORT FORM}\\
Memory variables & \cmd{STORE}, \cmd{RELEASE}, \cmd{PUBLIC}, \cmd{PRIVATE}\\
Questions to the user & \cmd{ACCEPT}, \cmd{INPUT}, \cmd{WAIT}\\
Programs & \cmd{DO}, \cmd{PROCEDURE}, \cmd{PARAMETERS}, \cmd{RETURN},
  \cmd{IF}, \cmd{DO CASE}, \cmd{DO WHILE}, \cmd{NOTE}, \cmd{ON ERROR}\\
The environment & the SET commands: SET CENTURY, SET DATE, SET DECIMALS, SET
  DELETED, SET EXACT, SET FILTER, SET HEADING, SET ORDER, SET SAFETY, SET
  STATUS and the others of Chapter~5\\
\hline
\end{tabular}

\section{The Sample Programs and Utilities Diskette}
The diskette holds the databases, index files, report forms and programs used
in the three sections of this manual: the staff roster \file{STAFF.DBF}, the
departments \file{DEPTS.DBF}, the accounts receivable ledger
\file{INVOICE.DBF}, their index files, the report forms \file{STAFFRPT.FRM} and
\file{INVRPT.FRM}, and the sample programs.  \doctitle{Getting Started} lists
every file on it.
